\section{Why defect matter reproduces the archimedean exponents}\label{sec:free}

The geometric reason for using three-dimensional defect matter can be
seen without supersymmetry. Normalize a free massless scalar covariance
on the defect to $|X-Y|^{-1}$; coupling and color constants will be
restored when needed. Write $X=e^x\Omega$, $Y=e^y\Omega'$ with
$\Omega,\Omega'\in S^2$, and rescale the field by $e^{x/2}$. With
$u=x-y$ and $v=\Omega\cdot\Omega'$, the covariance becomes
\begin{equation}\label{eq:cylinder-cov}
 C(u,v)=\frac1{\sqrt{2(\cosh u-v)}}
 =\sum_{\ell=0}^\infty e^{-(\ell+1/2)|u|}P_\ell(v),\qquad u\ne0.
\end{equation}
This is the Legendre generating function with parameter $e^{-|u|}$.

Fix an axis $e\in S^2$ and $0\leq\rho<1$. The Poisson density on $S^2$
and its even part are
\begin{equation}
 P_\rho(v)=\frac{1-\rho^2}{4\pi(1-2\rho v+\rho^2)^{3/2}},\qquad
 h_\rho(\Omega)=\tfrac12\{P_\rho(e\cdot\Omega)+P_\rho(-e\cdot\Omega)\}.
\end{equation}
Integrating \eqref{eq:cylinder-cov} against $h_\rho(\Omega)h_\rho(\Omega')$
projects onto even angular degrees. Orthogonality of the Legendre
polynomials gives
\begin{equation}\label{eq:regulated-cov}
 C_\rho(u)=\sum_{n=0}^\infty\rho^{4n}e^{-a_n|u|}
 =\frac{e^{-|u|/2}}{1-\rho^4e^{-2|u|}},\qquad a_n=2n+\tfrac12.
\end{equation}
For $\rho<1$ this is a bounded continuous positive-definite covariance.
For $u\ne0$ it tends to $\nGamma(|u|)$ as $\rho\uparrow1$.
The limit is singular at $u=0$, so positivity of the regulated covariance
cannot simply be passed to a finite contact-subtracted covariance.

\begin{proposition}[The sign of the subtraction]\label{prop:subtraction}
Let $T_\rho$ be convolution by $C_\rho$ on $\R$, and put
$M_\rho=\widehat C_\rho(0)$. Then $M_\rho I-T_\rho$ is positive,
and its quadratic form tends to $\mathcal E_\Gamma$ on smooth compactly
supported functions. The opposite subtraction $T_\rho-M_\rho I$ tends
to the negative of that form. Adding any finite constant multiple of
the identity to the latter limit does not give a positive form on all
smooth compactly supported functions.
\end{proposition}
\begin{proof}
Fourier transformation gives
\begin{align}
 \widehat C_\rho(\tau)&=\sum_{n\geq0}\rho^{4n}
 \frac{2a_n}{a_n^2+\tau^2},\notag\\
 M_\rho-\widehat C_\rho(\tau)
 &=2\sum_{n\geq0}\rho^{4n}\frac{\tau^2}{a_n(a_n^2+\tau^2)}
 \uparrow b(\tau^2).
\end{align}
Monotone convergence and \eqref{eq:kinetic} prove the first claims.
Moreover $M_\rho=\sum_{n\geq0}\rho^{4n}/(n+1/4)$ diverges
logarithmically. Since $b(\tau^2)\to\infty$, any multiplier
$c-b(\tau^2)$ is negative at high frequencies. Modulating a fixed
nonzero compactly supported smooth function to arbitrarily high
frequency gives a negative form value, also on any fixed nonempty
interval.
\end{proof}

There is a second useful decomposition. For $u\ne0$ define the
same-ray and opposite-ray functions
\begin{equation}\label{eq:rays}
 G_s(u)=\frac1{2\sinh(|u|/2)},\qquad
 G_o(u)=\frac1{2\cosh(u/2)}.
\end{equation}
Then $\nGamma(|u|)=\tfrac12(G_s(u)+G_o(u))$. The reflected free
network constructed in Section~\ref{sec:reflection} gives $G_o$.
That is a positive and finite kernel, but it is not the even tower by
itself. Conversely, the full angular smearing above is an ordinary
covariance construction with support on all of $S^2$. It is not
automatically a state supported in a positive-time half-space, as
required for a time-reflection argument. Restricting the angular
profile to a hemisphere changes its spectral weights.

The useful conclusion is therefore precise: free defect matter
supplies the desired exponents and a regulated positive covariance.
The RH target requires the associated difference energy, its fixed
local term, its pole term and its arithmetic translations. These
are distinct operations on the free data.

\subsection{An auxiliary Gaussian realization of the gamma ratio}

The parent endpoint work also gave a useful fixed-space alternative
to fractional dimension~\cite{ParentAngular}. It can be understood
from elementary Gaussian integration. Retain one complex mode for
each $a_n=2n+1/2$, and for real $p$ with $a_0+p>\omega$ define
\begin{equation}
 Z_N(p)=\prod_{n=0}^{N-1}\int_\C\frac{\mathrm d^2z_n}{\pi}
             e^{-(a_n+p)|z_n|^2}
       =\prod_{n=0}^{N-1}(a_n+p)^{-1}.
\end{equation}
For a half-line field $\phi_n(t)$ with action
$\int_0^\infty(|\phi_n'|^2+a_n^2|\phi_n|^2)\dd t$ and boundary value
$z_n$, the decaying classical solution is $\phi_n(t)=z_ne^{-a_nt}$.
Its action is $a_n|z_n|^2$. Adding the boundary term $p|z_n|^2$
gives the same Gaussian; varying the boundary value gives the Robin
condition $(-\partial_t+p)\phi_n(0)=0$ for the full variational problem.
Constants common to the two partition functions cancel in the ratio.

Set $x=(p+1/2-\omega)/2$ and $y=x+\omega$. The finite product gives
\begin{equation}
 \frac{Z_N(p-\omega)}{Z_N(p+\omega)}
 =\prod_{n=0}^{N-1}\frac{n+y}{n+x}
 =\frac{\Gamma(N+y)\Gamma(x)}{\Gamma(y)\Gamma(N+x)}.
\end{equation}
Since $\Gamma(N+y)/\Gamma(N+x)\sim N^{\omega}$,
\begin{equation}\label{eq:gaussian-gamma}
 \pi^\omega\lim_{N\to\infty}N^{-\omega}
 \frac{Z_N(p-\omega)}{Z_N(p+\omega)}
 =\pi^\omega\frac{\Gamma((p+a)/2)}{\Gamma((p+b)/2)}.
\end{equation}
Equivalently one can divide by the same finite product at a reference
$p_0$ to cancel the normalization. One real Gaussian mode would give
the square root of the ratio, so statistics and multiplicities matter.

This is an auxiliary partition-function model, not a localization
of the Wilson network. Angular smearing samples a channel; it does
not remove the other angular modes from a partition function. The
full scalar sphere has multiplicity $2\ell+1$, whereas the model
retains one copy of each even axial mode. A physical projection or
cancellation of the extra modes has not been derived. Here $p$ is a
boundary mass, and identifying it with the arithmetic Laplace
frequency requires a dynamical dictionary. The Robin response
$(a_n+p)^{-1}$, the partition ratio, and a causal transfer are different
objects. The supplied factor $\pi^\omega$, the prime comb, and the
rational pole correction remain outside this Gaussian mechanism.
